\section{模型评价、结论与建议}

\subsection{模型评价}

本文模型的优点主要有四点。第一，问题一把随机森林与 ALR 标定拆成“连续拟合”和“等级映射”两层，前者负责非线性与交互，后者负责题面要求的整数刻度，两层各自可检验，出了问题能定位到具体一层。第二，模型选择采用两阶段预注册，阈值、折索引、特征集都在结果产生之前冻结，避免了“看测试集挑模型”的选择偏差。第三，问题二把直接相关与图平滑写成可加融合分，并用 bootstrap、阈值扰动、负边消融与权重扰动四种方式检查，结论不依赖单一设置。第四，本文没有把 T2、T7、T8 的偏离包装成“模型很好”，而是给出刻度不可达命题与可读口径辅助表，既保留主表的契约口径，又让读者能按题面等级阅读。

模型的不足也有三点。第一，平均 $R^2$ 只有 \textbf{0.2305}（折均）与 \textbf{0.2565}（OOF 全局），部分弱维的逐维 $R^2$ 接近零甚至为负，说明化合物信息不足以解释所有感官差异。第二，问题二的 Top-1 较稳，但 Top-3 后两位的 bootstrap 频率偏低，第三名平均只有 \textbf{0.3611}，实际使用时不宜把 Top-3 当作同等可靠的结论。第三，问题三尚未开展实验，风格条件生成与跨风格迁移只有路线，没有结果；这是本阶段的边界，不是可以含糊带过的部分。

\subsection{结论}

\textbf{问题一}：随机森林与 \textbf{226} 维 F0 特征集是当前数据下的合理选择，OOF 全局 $R^2$ 为 \textbf{0.2565}，折均 $R^2$ 为 \textbf{0.2305}，均高于同口径基线。ALR 把连续预测转成 \textbf{0--6} 级；主表 \textbf{0} 级与 \textbf{6} 级为空是刻度不可达命题的必然结果，不是实现缺陷。可读口径辅助表在保持同序（逆序对 \textbf{0}）的前提下让 \textbf{7} 级全部出现，解决了“看起来没有 \textbf{6} 级”的阅读问题。

\textbf{问题二}：\textbf{226} 个候选化合物、\textbf{3758} 条边构成共表达网络，融合分在 $\lambda=0.70$ 下给出每个感官维度的 Top-3。Top-1 的平均入选频率为 \textbf{0.6734}，Top-3 为 \textbf{0.3611}；图项使 Top-3 相对直接相关的平均 Jaccard 降到 \textbf{0.378}，说明图结构确实参与排序。品类一致率为 \textbf{0.3582}，高于独立基线 \textbf{0.2047}，但离半数仍有距离；GBM 的 \textbf{60/60} 重合是恒真校验，有信息量的是平均分位 \textbf{0.6201} 对 \textbf{0.4971}。这些结果支持把 Top-3 作为实验候选集，而不是唯一因果解释。

\textbf{问题三}：待 M6 补。

\subsection{建议}

对酿造实践，本文给出三条建议。第一，优先用各维度的 Top-1 做靶向实验，把 Top-2 与 Top-3 作为备选；Top-1 的 bootstrap 频率明显高于后两位，直接拿第三名做工艺调整的风险较大。第二，使用评分表时同时保留主表与可读口径表：主表用于与训练标签同刻度的契约口径，可读表用于按题面等级展示；两者单调同序，不能互相替代，也不应把可读表当成新的训练目标。第三，后续数据采集应优先补强弱维与少数风格：弱维的逐维 $R^2$ 低，说明现有化合物面板缺少能解释这些感官的信息；少数风格样本少，会直接影响问题三的条件生成与迁移评估。

对后续建模，本文建议在问题三中把“化学有效性”放在“感官拟合”之前：先保证生成谱非负、量纲一致、关键化合物可追溯，再谈感官画像是否像目标风格。回评器可以使用问题一的随机森林，但必须对弱维降权，并用独立证据交叉核对，避免生成器利用回评器的误差“钻空子”。
